\documentclass[times]{san}

\usepackage{graphicx}
\usepackage{subcaption}
\usepackage{float}

\bibliografia{bibliografia.bib}

\uczelnia{Społeczna Akademia Nauk}
\przedmiot{Systemy Szkieletowe}
\autor{Maksym Mokriakov}
\nralbumu{121261}
\grupa{2}
\tytul{Debilingo -- aplikacja webowa wspierająca naukę języków obcych z elementami grywalizacji}
\prowadzacy{mgr inż. Krystian Gumiński}
\miasto{Łódź}
\rok{2026}

\begin{document}

\maketitlepage
\spistresci

\section{Wstęp}

Nauka języków obcych jest procesem długotrwałym, który wymaga regularności i dużego zaangażowania. W dobie powszechnego dostępu do urządzeń mobilnych, aplikacje wspomagające naukę stały się niezwykle popularne, jednak wiele z nich traci użytkowników z powodu szybkiego spadku motywacji po pierwszych dniach korzystania. 

Niniejszy projekt, noszący nazwę "Debilingo", ma na celu rozwiązanie tego problemu. Aplikacja łączy w sobie funkcjonalność osobistego słownika do nauki obcych słówek z silnymi mechanizmami grywalizacji (gamifikacji). Użytkownicy są nagradzani za regularność punktami doświadczenia (XP), rozwijają swój wirtualny poziom oraz budują "Daily streak" (codzienną serię). W raporcie przedstawiono architekturę rozwiązania, zastosowane technologie oraz kluczowe aspekty implementacyjne.

\section{Cel projektu}

Głównym celem projektu jest stworzenie nowoczesnej i responsywnej aplikacji webowej, zoptymalizowanej pod kątem urządzeń mobilnych, która umożliwi użytkownikom efektywną naukę słownictwa poprzez samodzielnie tworzone listy słów połączone z elementami motywującymi.

Wymagania funkcjonalne obejmują:
\begin{itemize}
    \item Rejestrację, logowanie oraz zarządzanie profilem użytkownika.
    \item Tworzenie własnych, spersonalizowanych słowników językowych.
    \item Dodawanie par słówek do wybranych słowników.
    \item Śledzenie postępów użytkownika (system poziomów, naliczanie punktów XP, codzienne serie).
\end{itemize}

Wymagania niefunkcjonalne to między innymi wysoka wydajność, niezawodność, zastosowanie architektury "Mobile-first" oraz pełne pokrycie kluczowych funkcjonalności testami automatycznymi. Skalowalność systemu zapewniono dzięki chmurze obliczeniowej i konteneryzacji.

\section{Przegląd istniejących rozwiązań}

Na rynku istnieje wiele aplikacji wspierających naukę języków, z których każda kładzie nacisk na inne aspekty edukacji. Analiza tych systemów pozwoliła na zidentyfikowanie zarówno dobrych praktyk, jak i obszarów wymagających poprawy.

\subsection{Duolingo}
Duolingo jest obecnie jedną z najpopularniejszych aplikacji do nauki języków na świecie. Jej ogromną zaletą jest niezwykle rozbudowany system grywalizacji, który silnie motywuje do codziennych ćwiczeń. Wadą dla zaawansowanych użytkowników jest jednak z góry narzucony program nauczania i brak możliwości wprowadzania własnego słownictwa z określonej branży (np. IT).

\subsection{Quizlet}
Quizlet to platforma oparta na wirtualnych fiszkach (flashcards), pozwalająca na tworzenie własnych zestawów słówek. Jest to bardzo elastyczne narzędzie, często wykorzystywane w szkołach i na uczelniach. Niestety, w darmowej wersji brakuje zaawansowanych mechanizmów podtrzymujących zaangażowanie i długotrwałą motywację (elementów gry), przez co proces nauki po pewnym czasie może wydawać się nużący.

\subsection{Anki}
Anki to potężne oprogramowanie wykorzystujące algorytm Spaced Repetition (rozłożonych w czasie powtórek). Pozwala na naukę czegokolwiek, od słówek po medycynę. Choć jest niesamowicie skuteczne, jego interfejs jest archaiczny i skomplikowany dla nowych użytkowników, a brak jakiejkolwiek grywalizacji sprawia, że wymaga żelaznej dyscypliny.

\vspace{0.5cm}
Debilingo stara się połączyć najlepsze cechy wyżej wymienionych systemów: elastyczność w dodawaniu własnych słówek (Quizlet) oraz wciągające mechanizmy motywacyjne i łatwy w obsłudze interfejs (Duolingo). Na rysunku \ref{fig:screen1} przedstawiono poglądowy widok aplikacji realizującej te założenia.



\section{Opis architektury}

System został zaprojektowany z wykorzystaniem architektury rozproszonej typu klient-serwer, opierając się o nowoczesne rozwiązania chmurowe. Tabela \ref{tab:architektura} prezentuje podział na główne warstwy.

\begin{table}[h]
\centering
\caption{Główne komponenty architektury systemu}
\label{tab:architektura}
\begin{tabular}{|c|p{4cm}|p{8cm}|}
\hline
Lp. & Komponent          & Technologia i rola \\
\hline
1   & Frontend (Aplikacja kliencka) & Vercel. Odpowiada za warstwę prezentacji (UI/UX). Oparta o React. \\
2   & Backend (API)      & Google Cloud Run. Realizuje logikę biznesową (Node.js, Express) z zachowaniem pełnej skalowalności. \\
3   & Baza danych        & Supabase (PostgreSQL). Platforma BaaS przechowująca dane użytkowników i słowniki. \\
\hline
\end{tabular}
\end{table}

Aplikacja frontendowa komunikuje się z serwerem backendowym za pośrednictwem żądań HTTP (REST API). Część zadań, takich jak natychmiastowe autoryzowanie sesji, wspierana jest również bezpośrednio przez klienta biblioteki Supabase w przeglądarce. Serwer backendowy jest całkowicie bezstanowy, co umożliwia jego bezpieczną konteneryzację w usłudze Google Cloud Run.

\section{Opis implementacji}

Implementacja systemu została podzielona na odrębne moduły zarządzane w ramach struktury monorepozytorium (oddzielne katalogi \texttt{web-app} oraz \texttt{backend}). 

\subsection{Frontend}
Frontend został zaimplementowany jako aplikacja Single Page Application (SPA) przy użyciu biblioteki React \cite{react} (w wersji 19) oraz języka TypeScript \cite{ts}. Jako narzędzie budujące wykorzystano Vite \cite{vite}, co znacząco przyspieszyło proces developmentu. Architektura kodu po stronie klienta oparta jest na nowoczesnej metodologii Feature-Sliced Design (FSD) \cite{fsd}, co widać w strukturze katalogów (\texttt{app}, \texttt{entities}, \texttt{features}, \texttt{pages}, \texttt{shared}, \texttt{widgets}). Zarządzanie stanem globalnym, w tym śledzenie punktów XP oraz trwającej serii użytkownika, powierzono bibliotece Zustand \cite{zustand}.

Interfejs został zaprojektowany zgodnie z podejściem Mobile-first.
% Interfejs został zaprojektowany zgodnie z podejściem Mobile-first (Rysunek \ref{fig:screen2}).

% \begin{figure}[h]
%     \centering
%     \includegraphics[width=0.45\textwidth]{image-1.png}
%     \caption{Poglądowy interfejs aplikacji mobilnej -- podsumowanie postępów}
%     \label{fig:screen1}
%     \end{figure}



    \begin{figure}[h]
        \centering
    
        \begin{subfigure}[t]{0.45\textwidth}
            \centering
            \includegraphics[height=0.55\textheight]{image-1.png}
            \caption{Strona główna}
            \label{fig:screen1}
        \end{subfigure}
        \hfill
        \begin{subfigure}[t]{0.45\textwidth}
            \centering
            \includegraphics[height=0.55\textheight]{image-2.png}
            \caption{Słownik}
            \label{fig:screen2}
        \end{subfigure}
    
        \caption{Poglądowy interfejs aplikacji mobilnej}
        \label{fig:two_screens}
    \end{figure}


\subsection{Backend}
Warstwa serwerowa została napisana z wykorzystaniem środowiska uruchomieniowego Node.js oraz frameworka Express.js. Całość zaimplementowano w języku TypeScript dla zachowania typowania (Type-Safety) na każdym etapie. Głównym zadaniem backendu jest obsługa zaawansowanej logiki biznesowej, bezpieczna wymiana danych oraz zarządzanie słownikami użytkowników. 



\subsection{Baza Danych}
Projekt wykorzystuje usługę Supabase \cite{supabase}, będącą otwartoźródłową alternatywą dla Firebase, pod spodem opierającą się na relacyjnej bazie danych PostgreSQL. Dzięki zastosowaniu mechanizmów Row Level Security (RLS) wbudowanych w Supabase, dane są całkowicie bezpieczne na poziomie bazy -- każdy użytkownik ma wgląd i możliwość modyfikacji wyłącznie własnych słowników i postępów. 


\subsection{Docker}
W celu łatwego i bezawaryjnego przenoszenia środowiska aplikacyjnego, backend został w pełni skonteneryzowany z użyciem technologii Docker \cite{docker}. Zbudowany obraz kontenera zawierający Node.js, zainstalowane zależności (\texttt{node\_modules}) i zbudowane pliki źródłowe, pozwala na uruchomienie aplikacji w dowolnym środowisku produkcyjnym -- ostatecznie wybrano elastyczną usługę Google Cloud Run \cite{gcr}, która skaluje ilość instancji w zależności od obciążenia.

\begin{figure}[H]
    \centering
    \includegraphics[width=0.5\textwidth]{dockerfile.png}
    \caption{Dockerfile}
    \label{fig:one_screen}
\end{figure}

\subsection{Testy}
Biorąc pod uwagę jakość wytworzonego oprogramowania, wprowadzono rygorystyczne procedury testowe. Do testów jednostkowych (Unit tests) na frontendzie użyto nowoczesnego i niezwykle szybkiego narzędzia Vitest, które współpracuje z systemem budowania Vite. Dodatkowo, w projekcie zaimplementowano testy end-to-end (E2E) korzystając z frameworka Playwright, które symulują realne zachowanie użytkownika przechodzącego przez ekrany aplikacji. Kod sprawdzany jest również linterem (ESLint) oraz formaterem (Prettier).

\section{Podsumowanie i wnioski}

W ramach zrealizowanego projektu stworzono w pełni funkcjonalną i nowoczesną aplikację wspomagającą naukę języków obcych (Debilingo). Udało się pomyślnie zintegrować mechanizmy grywalizacji z narzędziem typu "flashcards", co znacząco zwiększa motywację do codziennej nauki.
Zastosowanie technologii takich jak React 19, TypeScript oraz Supabase znacznie skróciło czas tworzenia oprogramowania i pozwoliło skupić się na dostarczeniu intuicyjnego interfejsu (UX).

Największym wyzwaniem w projekcie okazało się prawidłowe zaprojektowanie relacji w bazie danych PostgreSQL oraz obsługa stanu globalnego na frontendzie, co zostało skutecznie rozwiązane dzięki elastyczności biblioteki Zustand. System jest otwarty na rozbudowę. W przyszłości planowane jest dodanie m.in. systemu znajomych, globalnych rankingów użytkowników oraz generowania przykładowych zdań dla nowych słówek przy pomocy sztucznej inteligencji.

\spistabel
\spisrysunkow
\printbibliography

\end{document}